% Auto-generated LaTeX tables
% Include in paper with: % Auto-generated LaTeX tables
% Include in paper with: % Auto-generated LaTeX tables
% Include in paper with: % Auto-generated LaTeX tables
% Include in paper with: \input{tables/all_tables}



% ============================================================% Table: table1_main_results.tex
\begin{table}[t]
\centering
\caption{Evaluation results on MLLMU-Bench, UnLOK-VQA, and MMUBench (mean over 3 seeds). $\downarrow$ = lower is better, $\uparrow$ = higher is better. \textbf{Bold} = best per metric. UQS uses weights derived from retrained-model correlation (Table~\ref{tab:uqs_weights}).}
\label{tab:main_results}
\resizebox{\textwidth}{!}{%
\begin{tabular}{lcccccc}
\toprule
\textbf{Method} & FA$\downarrow$ & RA$\uparrow$ & MIA$\downarrow$ & AD$\downarrow$ & JS$\downarrow$ & UQS$\uparrow$ \\
\midrule
Gradient Ascent & 0.6222 & 0.6200 & \textbf{0.2822} & 106.1719 & 0.0654 & \textbf{0.6366} \\
Random Labels & \textbf{0.6111} & \textbf{0.6267} & 0.3378 & \textbf{41.6354} & \textbf{0.0393} & 0.6282 \\
FT-Retain & 0.6222 & 0.5800 & 0.3511 & 87.4271 & 0.1569 & 0.5899 \\
SalUn & 0.6222 & 0.6133 & 0.3222 & 50.9896 & 0.0409 & 0.6224 \\
\bottomrule
\end{tabular}%
}
\end{table}

% Table: table2_contradiction.tex
\begin{table}[t]
\centering
\caption{\textbf{The Contradiction.} Rank of each method under each metric (1 = best). Same method receives dramatically different ranks depending on which metric is used � e.g., Gradient Ascent is \#1 by FA but \#4 by RA. UQS provides a stable, principled ranking.}
\label{tab:contradiction}
\begin{tabular}{lcccccc}
\toprule
\textbf{Method} & \textbf{FA} & \textbf{RA} & \textbf{MIA} & \textbf{AD} & \textbf{JS} & \textbf{UQS} \\
\midrule
Gradient Ascent & 2 & 2 & \textbf{1} & 4 & 3 & \textbf{1} \\
Random Labels & \textbf{1} & \textbf{1} & 3 & \textbf{1} & \textbf{1} & 2 \\
FT-Retain & 3 & 4 & 4 & 3 & 4 & 4 \\
SalUn & 4 & 3 & 2 & 2 & 2 & 3 \\
\bottomrule
\end{tabular}
\end{table}

% Table: table3_tau_matrix.tex
\begin{table}[t]
\centering
\caption{Kendall's $\tau$ correlation between metric-induced method rankings (averaged over 3 datasets $\times$ 3 seeds). Near-zero or negative $\tau$ indicates metrics \textit{contradict} each other. \colorbox{gray!20}{Shaded} cells: $|\tau| < 0.3$ (contradiction region).}
\label{tab:tau_matrix}
\begin{tabular}{lccccc}
\toprule
 & \textbf{FA} & \textbf{RA} & \textbf{MIA} & \textbf{AD} & \textbf{JS} \\
\midrule
\textbf{FA} & � & +0.89 & +0.56 & \cellcolor{gray!20}\textbf{-0.22} & \cellcolor{gray!20}\textbf{-0.11} \\
\textbf{RA} & +0.89 & � & +0.67 & \cellcolor{gray!20}\textbf{-0.11} & \cellcolor{gray!20}\textbf{+0.00} \\
\textbf{MIA} & +0.56 & +0.67 & � & \cellcolor{gray!20}\textbf{-0.00} & \cellcolor{gray!20}\textbf{-0.11} \\
\textbf{AD} & \cellcolor{gray!20}\textbf{-0.22} & \cellcolor{gray!20}\textbf{-0.11} & \cellcolor{gray!20}\textbf{-0.00} & � & +0.44 \\
\textbf{JS} & \cellcolor{gray!20}\textbf{-0.11} & \cellcolor{gray!20}\textbf{+0.00} & \cellcolor{gray!20}\textbf{-0.11} & +0.44 & � \\
\bottomrule
\end{tabular}
\end{table}

% Table: table4_modality.tex
\begin{table}[t]
\centering
\caption{Mean pairwise Kendall's $\tau$ between evaluation metrics in multimodal (VQA) vs unimodal (CIFAR-10) settings. Lower $\tau$ in multimodal indicates greater metric disagreement, showing that the dual image+text pathway amplifies metric contradiction.}
\label{tab:modality}
\begin{tabular}{lcc}
\toprule
\textbf{Setting} & \textbf{Mean Pairwise $\tau$} & \textbf{$\Delta$ vs Unimodal} \\
\midrule
Multimodal (VQA) & \textbf{0.094} & -0.064 \\
Unimodal (CIFAR-10) & 0.158 & � \\
\bottomrule
\end{tabular}
\end{table}

% Table: table5_uqs_weights.tex
\begin{table}[t]
\centering
\caption{Metric reliability (Spearman $\rho$ with retrained-model distance) and derived UQS weights. Metrics more predictive of closeness to the gold-standard retrained model receive higher weight. All weights sum to 1.}
\label{tab:uqs_weights}
\begin{tabular}{lccl}
\toprule
\textbf{Metric} & \textbf{Spearman $\rho$} & \textbf{UQS Weight} & \textbf{Direction} \\
\midrule
\textbf{RA} & 0.482 & 0.705 & $\uparrow$ \\
\textbf{MIA} & 0.171 & 0.251 & $\downarrow$ \\
\textbf{FA} & -0.420 & 0.015 & $\downarrow$ \\
\textbf{AD} & -0.196 & 0.015 & $\downarrow$ \\
\textbf{JS} & -0.063 & 0.015 & $\downarrow$ \\
\midrule
Total & � & 1.000 & � \\
\bottomrule
\multicolumn{4}{l}{\small * $p<0.05$, ** $p<0.01$, *** $p<0.001$} \\
\end{tabular}
\end{table}

% Table: table6_uqs_ablation.tex
\begin{table}[t]
\centering
\caption{UQS ranking stability under random weight variations ($N=500$ trials, Dirichlet-sampled weights). High mean Kendall's $\tau$ shows that UQS rankings are stable regardless of exact weight values.}
\label{tab:uqs_ablation}
\begin{tabular}{lcc}
\toprule
\textbf{Metric} & \textbf{Value} & \textbf{Interpretation} \\
\midrule
Mean Kendall's $\tau$ & 0.599 & Rank stability \\
Std Kendall's $\tau$ & 0.265 & Variability \\
Number of trials & 500 & \\
\bottomrule
\end{tabular}
\end{table}



% ============================================================% Table: table1_main_results.tex
\begin{table}[t]
\centering
\caption{Evaluation results on MLLMU-Bench, UnLOK-VQA, and MMUBench (mean over 3 seeds). $\downarrow$ = lower is better, $\uparrow$ = higher is better. \textbf{Bold} = best per metric. UQS uses weights derived from retrained-model correlation (Table~\ref{tab:uqs_weights}).}
\label{tab:main_results}
\resizebox{\textwidth}{!}{%
\begin{tabular}{lcccccc}
\toprule
\textbf{Method} & FA$\downarrow$ & RA$\uparrow$ & MIA$\downarrow$ & AD$\downarrow$ & JS$\downarrow$ & UQS$\uparrow$ \\
\midrule
Gradient Ascent & 0.6222 & 0.6200 & \textbf{0.2822} & 106.1719 & 0.0654 & \textbf{0.6366} \\
Random Labels & \textbf{0.6111} & \textbf{0.6267} & 0.3378 & \textbf{41.6354} & \textbf{0.0393} & 0.6282 \\
FT-Retain & 0.6222 & 0.5800 & 0.3511 & 87.4271 & 0.1569 & 0.5899 \\
SalUn & 0.6222 & 0.6133 & 0.3222 & 50.9896 & 0.0409 & 0.6224 \\
\bottomrule
\end{tabular}%
}
\end{table}

% Table: table2_contradiction.tex
\begin{table}[t]
\centering
\caption{\textbf{The Contradiction.} Rank of each method under each metric (1 = best). Same method receives dramatically different ranks depending on which metric is used � e.g., Gradient Ascent is \#1 by FA but \#4 by RA. UQS provides a stable, principled ranking.}
\label{tab:contradiction}
\begin{tabular}{lcccccc}
\toprule
\textbf{Method} & \textbf{FA} & \textbf{RA} & \textbf{MIA} & \textbf{AD} & \textbf{JS} & \textbf{UQS} \\
\midrule
Gradient Ascent & 2 & 2 & \textbf{1} & 4 & 3 & \textbf{1} \\
Random Labels & \textbf{1} & \textbf{1} & 3 & \textbf{1} & \textbf{1} & 2 \\
FT-Retain & 3 & 4 & 4 & 3 & 4 & 4 \\
SalUn & 4 & 3 & 2 & 2 & 2 & 3 \\
\bottomrule
\end{tabular}
\end{table}

% Table: table3_tau_matrix.tex
\begin{table}[t]
\centering
\caption{Kendall's $\tau$ correlation between metric-induced method rankings (averaged over 3 datasets $\times$ 3 seeds). Near-zero or negative $\tau$ indicates metrics \textit{contradict} each other. \colorbox{gray!20}{Shaded} cells: $|\tau| < 0.3$ (contradiction region).}
\label{tab:tau_matrix}
\begin{tabular}{lccccc}
\toprule
 & \textbf{FA} & \textbf{RA} & \textbf{MIA} & \textbf{AD} & \textbf{JS} \\
\midrule
\textbf{FA} & � & +0.89 & +0.56 & \cellcolor{gray!20}\textbf{-0.22} & \cellcolor{gray!20}\textbf{-0.11} \\
\textbf{RA} & +0.89 & � & +0.67 & \cellcolor{gray!20}\textbf{-0.11} & \cellcolor{gray!20}\textbf{+0.00} \\
\textbf{MIA} & +0.56 & +0.67 & � & \cellcolor{gray!20}\textbf{-0.00} & \cellcolor{gray!20}\textbf{-0.11} \\
\textbf{AD} & \cellcolor{gray!20}\textbf{-0.22} & \cellcolor{gray!20}\textbf{-0.11} & \cellcolor{gray!20}\textbf{-0.00} & � & +0.44 \\
\textbf{JS} & \cellcolor{gray!20}\textbf{-0.11} & \cellcolor{gray!20}\textbf{+0.00} & \cellcolor{gray!20}\textbf{-0.11} & +0.44 & � \\
\bottomrule
\end{tabular}
\end{table}

% Table: table4_modality.tex
\begin{table}[t]
\centering
\caption{Mean pairwise Kendall's $\tau$ between evaluation metrics in multimodal (VQA) vs unimodal (CIFAR-10) settings. Lower $\tau$ in multimodal indicates greater metric disagreement, showing that the dual image+text pathway amplifies metric contradiction.}
\label{tab:modality}
\begin{tabular}{lcc}
\toprule
\textbf{Setting} & \textbf{Mean Pairwise $\tau$} & \textbf{$\Delta$ vs Unimodal} \\
\midrule
Multimodal (VQA) & \textbf{0.094} & -0.064 \\
Unimodal (CIFAR-10) & 0.158 & � \\
\bottomrule
\end{tabular}
\end{table}

% Table: table5_uqs_weights.tex
\begin{table}[t]
\centering
\caption{Metric reliability (Spearman $\rho$ with retrained-model distance) and derived UQS weights. Metrics more predictive of closeness to the gold-standard retrained model receive higher weight. All weights sum to 1.}
\label{tab:uqs_weights}
\begin{tabular}{lccl}
\toprule
\textbf{Metric} & \textbf{Spearman $\rho$} & \textbf{UQS Weight} & \textbf{Direction} \\
\midrule
\textbf{RA} & 0.482 & 0.705 & $\uparrow$ \\
\textbf{MIA} & 0.171 & 0.251 & $\downarrow$ \\
\textbf{FA} & -0.420 & 0.015 & $\downarrow$ \\
\textbf{AD} & -0.196 & 0.015 & $\downarrow$ \\
\textbf{JS} & -0.063 & 0.015 & $\downarrow$ \\
\midrule
Total & � & 1.000 & � \\
\bottomrule
\multicolumn{4}{l}{\small * $p<0.05$, ** $p<0.01$, *** $p<0.001$} \\
\end{tabular}
\end{table}

% Table: table6_uqs_ablation.tex
\begin{table}[t]
\centering
\caption{UQS ranking stability under random weight variations ($N=500$ trials, Dirichlet-sampled weights). High mean Kendall's $\tau$ shows that UQS rankings are stable regardless of exact weight values.}
\label{tab:uqs_ablation}
\begin{tabular}{lcc}
\toprule
\textbf{Metric} & \textbf{Value} & \textbf{Interpretation} \\
\midrule
Mean Kendall's $\tau$ & 0.599 & Rank stability \\
Std Kendall's $\tau$ & 0.265 & Variability \\
Number of trials & 500 & \\
\bottomrule
\end{tabular}
\end{table}



% ============================================================% Table: table1_main_results.tex
\begin{table}[t]
\centering
\caption{Evaluation results on MLLMU-Bench, UnLOK-VQA, and MMUBench (mean over 3 seeds). $\downarrow$ = lower is better, $\uparrow$ = higher is better. \textbf{Bold} = best per metric. UQS uses weights derived from retrained-model correlation (Table~\ref{tab:uqs_weights}).}
\label{tab:main_results}
\resizebox{\textwidth}{!}{%
\begin{tabular}{lcccccc}
\toprule
\textbf{Method} & FA$\downarrow$ & RA$\uparrow$ & MIA$\downarrow$ & AD$\downarrow$ & JS$\downarrow$ & UQS$\uparrow$ \\
\midrule
Gradient Ascent & 0.6222 & 0.6200 & \textbf{0.2822} & 106.1719 & 0.0654 & \textbf{0.6366} \\
Random Labels & \textbf{0.6111} & \textbf{0.6267} & 0.3378 & \textbf{41.6354} & \textbf{0.0393} & 0.6282 \\
FT-Retain & 0.6222 & 0.5800 & 0.3511 & 87.4271 & 0.1569 & 0.5899 \\
SalUn & 0.6222 & 0.6133 & 0.3222 & 50.9896 & 0.0409 & 0.6224 \\
\bottomrule
\end{tabular}%
}
\end{table}

% Table: table2_contradiction.tex
\begin{table}[t]
\centering
\caption{\textbf{The Contradiction.} Rank of each method under each metric (1 = best). Same method receives dramatically different ranks depending on which metric is used � e.g., Gradient Ascent is \#1 by FA but \#4 by RA. UQS provides a stable, principled ranking.}
\label{tab:contradiction}
\begin{tabular}{lcccccc}
\toprule
\textbf{Method} & \textbf{FA} & \textbf{RA} & \textbf{MIA} & \textbf{AD} & \textbf{JS} & \textbf{UQS} \\
\midrule
Gradient Ascent & 2 & 2 & \textbf{1} & 4 & 3 & \textbf{1} \\
Random Labels & \textbf{1} & \textbf{1} & 3 & \textbf{1} & \textbf{1} & 2 \\
FT-Retain & 3 & 4 & 4 & 3 & 4 & 4 \\
SalUn & 4 & 3 & 2 & 2 & 2 & 3 \\
\bottomrule
\end{tabular}
\end{table}

% Table: table3_tau_matrix.tex
\begin{table}[t]
\centering
\caption{Kendall's $\tau$ correlation between metric-induced method rankings (averaged over 3 datasets $\times$ 3 seeds). Near-zero or negative $\tau$ indicates metrics \textit{contradict} each other. \colorbox{gray!20}{Shaded} cells: $|\tau| < 0.3$ (contradiction region).}
\label{tab:tau_matrix}
\begin{tabular}{lccccc}
\toprule
 & \textbf{FA} & \textbf{RA} & \textbf{MIA} & \textbf{AD} & \textbf{JS} \\
\midrule
\textbf{FA} & � & +0.89 & +0.56 & \cellcolor{gray!20}\textbf{-0.22} & \cellcolor{gray!20}\textbf{-0.11} \\
\textbf{RA} & +0.89 & � & +0.67 & \cellcolor{gray!20}\textbf{-0.11} & \cellcolor{gray!20}\textbf{+0.00} \\
\textbf{MIA} & +0.56 & +0.67 & � & \cellcolor{gray!20}\textbf{-0.00} & \cellcolor{gray!20}\textbf{-0.11} \\
\textbf{AD} & \cellcolor{gray!20}\textbf{-0.22} & \cellcolor{gray!20}\textbf{-0.11} & \cellcolor{gray!20}\textbf{-0.00} & � & +0.44 \\
\textbf{JS} & \cellcolor{gray!20}\textbf{-0.11} & \cellcolor{gray!20}\textbf{+0.00} & \cellcolor{gray!20}\textbf{-0.11} & +0.44 & � \\
\bottomrule
\end{tabular}
\end{table}

% Table: table4_modality.tex
\begin{table}[t]
\centering
\caption{Mean pairwise Kendall's $\tau$ between evaluation metrics in multimodal (VQA) vs unimodal (CIFAR-10) settings. Lower $\tau$ in multimodal indicates greater metric disagreement, showing that the dual image+text pathway amplifies metric contradiction.}
\label{tab:modality}
\begin{tabular}{lcc}
\toprule
\textbf{Setting} & \textbf{Mean Pairwise $\tau$} & \textbf{$\Delta$ vs Unimodal} \\
\midrule
Multimodal (VQA) & \textbf{0.094} & -0.064 \\
Unimodal (CIFAR-10) & 0.158 & � \\
\bottomrule
\end{tabular}
\end{table}

% Table: table5_uqs_weights.tex
\begin{table}[t]
\centering
\caption{Metric reliability (Spearman $\rho$ with retrained-model distance) and derived UQS weights. Metrics more predictive of closeness to the gold-standard retrained model receive higher weight. All weights sum to 1.}
\label{tab:uqs_weights}
\begin{tabular}{lccl}
\toprule
\textbf{Metric} & \textbf{Spearman $\rho$} & \textbf{UQS Weight} & \textbf{Direction} \\
\midrule
\textbf{RA} & 0.482 & 0.705 & $\uparrow$ \\
\textbf{MIA} & 0.171 & 0.251 & $\downarrow$ \\
\textbf{FA} & -0.420 & 0.015 & $\downarrow$ \\
\textbf{AD} & -0.196 & 0.015 & $\downarrow$ \\
\textbf{JS} & -0.063 & 0.015 & $\downarrow$ \\
\midrule
Total & � & 1.000 & � \\
\bottomrule
\multicolumn{4}{l}{\small * $p<0.05$, ** $p<0.01$, *** $p<0.001$} \\
\end{tabular}
\end{table}

% Table: table6_uqs_ablation.tex
\begin{table}[t]
\centering
\caption{UQS ranking stability under random weight variations ($N=500$ trials, Dirichlet-sampled weights). High mean Kendall's $\tau$ shows that UQS rankings are stable regardless of exact weight values.}
\label{tab:uqs_ablation}
\begin{tabular}{lcc}
\toprule
\textbf{Metric} & \textbf{Value} & \textbf{Interpretation} \\
\midrule
Mean Kendall's $\tau$ & 0.599 & Rank stability \\
Std Kendall's $\tau$ & 0.265 & Variability \\
Number of trials & 500 & \\
\bottomrule
\end{tabular}
\end{table}



% ============================================================% Table: table1_main_results.tex
\begin{table}[t]
\centering
\caption{Evaluation results on MLLMU-Bench, UnLOK-VQA, and MMUBench (mean over 3 seeds). $\downarrow$ = lower is better, $\uparrow$ = higher is better. \textbf{Bold} = best per metric. UQS uses weights derived from retrained-model correlation (Table~\ref{tab:uqs_weights}).}
\label{tab:main_results}
\resizebox{\textwidth}{!}{%
\begin{tabular}{lcccccc}
\toprule
\textbf{Method} & FA$\downarrow$ & RA$\uparrow$ & MIA$\downarrow$ & AD$\downarrow$ & JS$\downarrow$ & UQS$\uparrow$ \\
\midrule
Gradient Ascent & 0.6222 & 0.6200 & \textbf{0.2822} & 106.1719 & 0.0654 & \textbf{0.6366} \\
Random Labels & \textbf{0.6111} & \textbf{0.6267} & 0.3378 & \textbf{41.6354} & \textbf{0.0393} & 0.6282 \\
FT-Retain & 0.6222 & 0.5800 & 0.3511 & 87.4271 & 0.1569 & 0.5899 \\
SalUn & 0.6222 & 0.6133 & 0.3222 & 50.9896 & 0.0409 & 0.6224 \\
\bottomrule
\end{tabular}%
}
\end{table}

% Table: table2_contradiction.tex
\begin{table}[t]
\centering
\caption{\textbf{The Contradiction.} Rank of each method under each metric (1 = best). Same method receives dramatically different ranks depending on which metric is used � e.g., Gradient Ascent is \#1 by FA but \#4 by RA. UQS provides a stable, principled ranking.}
\label{tab:contradiction}
\begin{tabular}{lcccccc}
\toprule
\textbf{Method} & \textbf{FA} & \textbf{RA} & \textbf{MIA} & \textbf{AD} & \textbf{JS} & \textbf{UQS} \\
\midrule
Gradient Ascent & 2 & 2 & \textbf{1} & 4 & 3 & \textbf{1} \\
Random Labels & \textbf{1} & \textbf{1} & 3 & \textbf{1} & \textbf{1} & 2 \\
FT-Retain & 3 & 4 & 4 & 3 & 4 & 4 \\
SalUn & 4 & 3 & 2 & 2 & 2 & 3 \\
\bottomrule
\end{tabular}
\end{table}

% Table: table3_tau_matrix.tex
\begin{table}[t]
\centering
\caption{Kendall's $\tau$ correlation between metric-induced method rankings (averaged over 3 datasets $\times$ 3 seeds). Near-zero or negative $\tau$ indicates metrics \textit{contradict} each other. \colorbox{gray!20}{Shaded} cells: $|\tau| < 0.3$ (contradiction region).}
\label{tab:tau_matrix}
\begin{tabular}{lccccc}
\toprule
 & \textbf{FA} & \textbf{RA} & \textbf{MIA} & \textbf{AD} & \textbf{JS} \\
\midrule
\textbf{FA} & � & +0.89 & +0.56 & \cellcolor{gray!20}\textbf{-0.22} & \cellcolor{gray!20}\textbf{-0.11} \\
\textbf{RA} & +0.89 & � & +0.67 & \cellcolor{gray!20}\textbf{-0.11} & \cellcolor{gray!20}\textbf{+0.00} \\
\textbf{MIA} & +0.56 & +0.67 & � & \cellcolor{gray!20}\textbf{-0.00} & \cellcolor{gray!20}\textbf{-0.11} \\
\textbf{AD} & \cellcolor{gray!20}\textbf{-0.22} & \cellcolor{gray!20}\textbf{-0.11} & \cellcolor{gray!20}\textbf{-0.00} & � & +0.44 \\
\textbf{JS} & \cellcolor{gray!20}\textbf{-0.11} & \cellcolor{gray!20}\textbf{+0.00} & \cellcolor{gray!20}\textbf{-0.11} & +0.44 & � \\
\bottomrule
\end{tabular}
\end{table}

% Table: table4_modality.tex
\begin{table}[t]
\centering
\caption{Mean pairwise Kendall's $\tau$ between evaluation metrics in multimodal (VQA) vs unimodal (CIFAR-10) settings. Lower $\tau$ in multimodal indicates greater metric disagreement, showing that the dual image+text pathway amplifies metric contradiction.}
\label{tab:modality}
\begin{tabular}{lcc}
\toprule
\textbf{Setting} & \textbf{Mean Pairwise $\tau$} & \textbf{$\Delta$ vs Unimodal} \\
\midrule
Multimodal (VQA) & \textbf{0.094} & -0.064 \\
Unimodal (CIFAR-10) & 0.158 & � \\
\bottomrule
\end{tabular}
\end{table}

% Table: table5_uqs_weights.tex
\begin{table}[t]
\centering
\caption{Metric reliability (Spearman $\rho$ with retrained-model distance) and derived UQS weights. Metrics more predictive of closeness to the gold-standard retrained model receive higher weight. All weights sum to 1.}
\label{tab:uqs_weights}
\begin{tabular}{lccl}
\toprule
\textbf{Metric} & \textbf{Spearman $\rho$} & \textbf{UQS Weight} & \textbf{Direction} \\
\midrule
\textbf{RA} & 0.482 & 0.705 & $\uparrow$ \\
\textbf{MIA} & 0.171 & 0.251 & $\downarrow$ \\
\textbf{FA} & -0.420 & 0.015 & $\downarrow$ \\
\textbf{AD} & -0.196 & 0.015 & $\downarrow$ \\
\textbf{JS} & -0.063 & 0.015 & $\downarrow$ \\
\midrule
Total & � & 1.000 & � \\
\bottomrule
\multicolumn{4}{l}{\small * $p<0.05$, ** $p<0.01$, *** $p<0.001$} \\
\end{tabular}
\end{table}

% Table: table6_uqs_ablation.tex
\begin{table}[t]
\centering
\caption{UQS ranking stability under random weight variations ($N=500$ trials, Dirichlet-sampled weights). High mean Kendall's $\tau$ shows that UQS rankings are stable regardless of exact weight values.}
\label{tab:uqs_ablation}
\begin{tabular}{lcc}
\toprule
\textbf{Metric} & \textbf{Value} & \textbf{Interpretation} \\
\midrule
Mean Kendall's $\tau$ & 0.599 & Rank stability \\
Std Kendall's $\tau$ & 0.265 & Variability \\
Number of trials & 500 & \\
\bottomrule
\end{tabular}
\end{table}